\documentclass[../../../include/open-logic-section]{subfiles}

\begin{document}

\olfileid{sth}{choice}{tarskiscott}
\olsection{حيلة تارسكي--سكوت}

عرّفنا في \olref[cardinals][cardsasords]{defcardinalasordinal}
الكاردينالات بوصفها أعداداً ترتيبية. ولفعل ذلك افترضنا بديهية الترتيب
الجيد. ولم نفعل ذلك لسبب سوى أنه «المعيار المتبع».

فقبل أن نناقش أياً من المسائل الفلسفية المحيطة بالترتيب الجيد، من المهم
أن يتضح أننا \emph{نستطيع} الخروج على المعيار المتبع وتطوير نظرية
للكاردينالات \emph{من دون} افتراض الترتيب الجيد. وما زال في وسعنا
استعمال تعريفات $\cardeq{A}{B}$ و$\cardle{A}{B}$ و$\cardless{A}{B}$
كما وردت في \olref[sfr][siz][]{chap}. ولن نحتاج إلا إلى مفهوم جديد
للـ\emph{كاردينال}.

وقد يخطر لنا، على نحو ساذج، أن نحاول تعريف كاردينالية $A$ هكذا:
\[
	\Setabs{x}{\cardeq{A}{x}}.
\]
ولعلك تقارن ذلك بتعريف فرغه لـ$\# x Fx$، الذي عرضنا خطوطه العامة في
نهاية \olref[cardinals][hp]{sec}. لكن هذا التعريف يفشل للأسباب التي
ألمحنا إليها هناك. خذ تحديداً $A=\{\emptyset\}$. فكل مجموعة أحادية
العنصر متساوية القدرة مع $A$، غير أن مجموعات أحادية جديدة تتكون عند كل
مرحلة خلفية من التراتبية (انظر فقط إلى المجموعة الأحادية للمرحلة
السابقة). ولذلك لا توجد $\Setabs{x}{\cardeq{A}{x}}$، إذ لا يمكن أن تكون
لها رتبة.

وللالتفاف على هذه المشكلة نستعمل حيلة تعزى إلى تارسكي
وسكوت:\footnote{تذكير: يجوز أن تتضمن جميع الصيغ وسائط، ما لم ينص صراحة
على خلاف ذلك.}

\begin{defn}[Tarski--Scott]
لكل صيغة $\phi(x)$، لتكن $[ x : \phi(x)] $ مجموعة جميع قيم $x$ ذات
أصغر رتبة ممكنة التي تحقق $\phi(x)$ (أو $\emptyset$ إذا لم توجد قيم
تحقق $\phi$).
\end{defn}

ينبغي أن نتحقق من مشروعية هذا التعريف. بالعمل في $\ZF$، تضمن
\olref[spine][foundation]{zfentailsregularity} وجود $\setrank{x}$ لكل
$x$. فإذا وجدت أشياء تحقق $\phi$، اخترنا أصغر $\alpha$ يحقق
$(\exists x\subseteq V_\alpha)\phi(x)$. وبحسب أصغرية $\alpha$ لا توجد
في أي رتبة أدنى قيمة تحقق $\phi$؛ أي إن
$(\forall \beta \in \alpha)(\forall x \subseteq V_\beta)
\lnot \phi(x)$. ويمكننا عندئذ، باستعمال الفصل، أن نعرّف
$[x : \phi(x)]$ بأنه $\Setabs{x \in V_{\alpha+1}}{\phi(x)}$.

وبعد أن بررنا حيلة تارسكي--سكوت، يمكننا استعمالها لتعريف مفهوم
للكاردينالية:

\begin{defn}
كاردينالية \textsc{ts} لـ$A$ هي
$\text{tsc}(A) = [x : \cardeq{A}{x}]$.
\end{defn}

لا يستعمل تعريف كاردينال \textsc{ts} الترتيب الجيد. لكننا نستطيع، حتى
من دون تلك البديهية، أن نبين أن \emph{كاردينالات \textsc{ts}} تتصرف إلى
حد كبير مثل \emph{الكاردينالات} كما عُرِّفت في
\olref[cardinals][cardsasords]{defcardinalasordinal}. فعلى سبيل المثال،
إذا أعدنا صياغة
\olref[cardinals][cardsasords]{lem:CardinalsBehaveRight}
و\olref[card-arithmetic][opps]{lem:SizePowerset2Exp} بلغة كاردينالات
\textsc{ts}، فإن البراهين تظل صحيحة في $\ZF$، من دون افتراض الترتيب
الجيد.

وبمناسبة هذا الموضوع، تجدر الإشارة إلى أننا نستطيع أيضاً تطوير نظرية
للأعداد الترتيبية باستعمال حيلة تارسكي--سكوت. فإذا كان
$\tuple{A, <}$ ترتيباً جيداً، فلتكن
$\text{tso}(A, <) = [\tuple{X, R} :
\ordeq{\tuple{A, <}}{\tuple{X, R}}]$. ولمزيد من التفصيل عن هذه المعالجة
للكاردينالات والأعداد الترتيبية، انظر
\citet[chs.~9--12]{Potter2004}.

\end{document}
